\section{行为主体 - 自然人}
\begin{enumerate}
  \item 自然人。是个人就行，\zd{有作案能力}即可，不要求责任年龄、责任能力。
  \item 身份犯。
        \begin{enumerate}
          \item 真正身份犯/定罪身份犯，只有该身份才能构成该罪。\al{贪污罪}。
                \begin{enumerate}
                  \item 开始\zd{犯罪时就要具备}；
                  \item 只针对该罪的实行犯。帮助犯/教唆犯构成不要求有这个身份。\al{妻子教唆丈夫贪污}。
                \end{enumerate}
          \item 不真正身份犯/量刑身份犯。谁都能构成，但某类构成，量刑上从重处罚。\al{诬告陷害罪}。
        \end{enumerate}
  \item \zd{国家工作人员}。
        \begin{enumerate}
          \item 国家工作人员（国有企、事业单位）|贪污罪 > 国家机关工作人员|渎职罪 > 司法机关工作人员|徇私枉法类。
          \item 身份认定。
                \begin{enumerate}
                  \item 从事\zd{公务}。公共管理性、行政职责性。
                  \item \zd{不看有没有编制}。区别于行政法。\al{协助政府从事公务，协警、辅警、疫情发物资贪污}。
                \end{enumerate}
          \item 身份具有切换性。国家工作人员/非国家工作人员。
                \begin{enumerate}
                  \item 贪钱。贪污/职务侵占。
                  \item 收钱。受贿/非国家工作人员受贿罪。
                  \item 挪用。挪用公款罪/挪用资金罪。
                \end{enumerate}
        \end{enumerate}
\end{enumerate}
